\section{Related work}
\label{sec:related}
\paragraph{Referring and reasoning segmentation.}
Referring image segmentation uses language to specify the object whose pixels
should be predicted. Cross-Modal Self-Attention Network~\cite{ye2019cross}
models interactions between linguistic and visual features for this task.
LISA~\cite{lai2024lisa} extends the instruction interface to reasoning-oriented
requests and decodes a segmentation-token representation into a mask.
These formulations establish that identifying the requested target can require
more than recognizing its category. Our question concerns an additional
conditioning variable: when a localized worker reference changes, does a
fixed relational request select the corresponding equipment instance?
The output remains a segmentation mask, but correctness includes following
the supplied identity intervention. Language understanding and mask quality
are therefore necessary without being sufficient evidence of identity control.
CMA inherits the LISA-style language-to-mask path and the promptable SAM
backbone~\cite{kirillov2023segment}; these inherited capabilities are not
presented as contributions of our reference-conditioning mechanism.

\paragraph{Visual reference and in-context segmentation.}
SAM~\cite{kirillov2023segment} provides a promptable interface for selecting
image regions. Visual references can additionally specify what should be
segmented through examples. VRP-SAM~\cite{sun2024vrp} encodes an annotated
reference image to guide segmentation in a target image, supporting several
forms of reference annotation. Visual In-Context Prompting, implemented by
DINOv~\cite{li2024visual}, supports referring and generic segmentation using
visual prompts and reference image segments, including single-image and
cross-image settings. These works establish relevant ways to communicate
visual intent; they should not be reduced to text-only segmentation methods.
For CMA, the supplied reference localizes a miner, whereas the requested
output is the helmet associated with that miner. The memory contains a miner
mask, bounding box and appearance crop from the same condition image.
Consequently, the evaluated task does not require discovering the worker
from an unlocalized identity token. Nor does it test cross-image appearance
transfer from a separately preserved clean exemplar. Our contrast with visual
reference segmentation concerns this particular input-to-target relation and
the controlled evaluation, rather than a claim that those architectures cannot
express relations. VRP-SAM and DINOv provide literature context here; neither
has been evaluated as a baseline in the frozen comparison.

\paragraph{Multi-round and memory-conditioned segmentation.}
SegLLM~\cite{wang2025segllm} explicitly reuses visual and textual outputs from
previous rounds and segments objects related to previously identified
entities. Its relational task therefore already overlaps the broad idea of
using entity context to select a different target. We do not claim priority
for that capability. CMA studies a narrower intervention: supplied localized
miner memory selects a relationally associated helmet, and switching that
memory while holding scene and query fixed tests whether target identity
follows the reference. The paired overlap matrix separates correct-identity
preference, successful localization of both identities, and confident
wrong-identity selection. The inspected literature does not establish an
identical evaluation contract; this observation is not an exhaustive absence
claim. Our pinned SegLLM comparison measures complete released/trained
systems through their native memory interfaces. It does not control training
exposure or isolate architectural effects. The resulting evidence is limited
to supplied-memory use on the frozen pseudo-labelled subset under one
compound stressor. Conversational memory in prior work also does not turn
our fixed-forward evaluation into autonomous memory acquisition, persistent
tracking, or an agent that repairs unreliable state.
